% ============================================================
% PORTADA
% ============================================================
{
  \setbeamertemplate{navigation symbols}{}
  \setbeamertemplate{footline}{}
  \begin{frame}[plain]
    \titlepage
  \end{frame}
}

% ============================================================
% TABLA DE CONTENIDOS
% ============================================================
{
  \setbeamertemplate{navigation symbols}{}
  \setbeamertemplate{footline}{}
  \begin{frame}{Outline}
    \tableofcontents
  \end{frame}
}

% ============================================================
\section{Problem Definition}
% ============================================================

\begin{frame}{What problem are we solving?}

  The open-source econometrics ecosystem is vast and fragmented.
  Researchers, central banks, and consulting firms need to answer:

  \bigskip
  \begin{block}{Central question}
    \centering
    \textit{Which quantitative economics tools are \textbf{gaining momentum},
    which are \textbf{established}, and which are \textbf{fading away}?}
  \end{block}

  \bigskip
  \begin{columns}
    \column{0.5\textwidth}
    \textbf{The challenge:}
    \begin{itemize}
      \item Thousands of GitHub repos — impossible to review manually
      \item No ground-truth labels for technology maturity
      \item Signals are noisy and multidimensional
    \end{itemize}

    \column{0.5\textwidth}
    \textbf{Our solution:}
    \begin{itemize}
      \item Collect 12 GitHub signals per repository
      \item Use DeepSeek (LLM) to generate \textit{weak labels}
      \item Fine-tune DistilBERT to learn the classification
    \end{itemize}
  \end{columns}

\end{frame}

% ------------------------------------------------------------

\begin{frame}{Four categories of technological maturity}

  \begin{columns}
    \column{0.48\textwidth}

    \begin{block}{\textcolor{emerging}{\textbf{Emerging}}}
      Fast-growing, recent momentum, high activity.
      \textit{Ex: new causal inference libraries, nowcasting tools.}
    \end{block}

    \vspace{0.4em}

    \begin{block}{\textcolor{mature}{\textbf{Mature}}}
      Stable, widely adopted, long track record.
      \textit{Ex: \texttt{statsmodels}, \texttt{linearmodels}.}
    \end{block}

    \column{0.48\textwidth}

    \begin{block}{\textcolor{declining}{\textbf{Declining}}}
      Decreasing activity, long inactivity periods.
      \textit{Ex: abandoned DSGE wrappers, unmaintained R ports.}
    \end{block}

    \vspace{0.4em}

    \begin{block}{\textcolor{experimental}{\textbf{Experimental / Niche}}}
      Early-stage or highly specialized, limited adoption.
      \textit{Ex: structural estimation solvers, VAR extensions.}
    \end{block}

  \end{columns}

  \bigskip
  \small These categories were designed to reflect \textbf{technological momentum},
  not code quality.

\end{frame}

% ============================================================
\section{Data Collection \& Signals}
% ============================================================

\begin{frame}{GitHub data collection}

  \begin{bluebox}{Collection strategy}
    587 repositories collected via 15 keyword/topic queries targeting the
    \textbf{Econometrics \& Macro Tooling} ecosystem using the GitHub REST API.
    Queries include: \texttt{econometrics}, \texttt{macroeconomics}, \texttt{dsge},
    \texttt{panel-data}, \texttt{time-series}, \texttt{cointegration}, \texttt{nowcasting},
    \texttt{causal-inference}, \texttt{structural-estimation}.
  \end{bluebox}

  \bigskip
  \textbf{Why this ecosystem?}
  \begin{itemize}
    \item Natural mix of the four maturity categories (mature tools like
          \texttt{statsmodels} coexist with experimental DSGE solvers)
    \item Domain expertise enables better justification of design choices
    \item Business relevance: central banks, IMF, consulting firms rely on these tools
  \end{itemize}

  \bigskip
  \small \textbf{Potential bias:} GitHub search favors popular, well-maintained repos.
  Niche or abandoned tools may be underrepresented.

\end{frame}

% ------------------------------------------------------------

\begin{frame}{12 repository signals}

  \begin{columns}
    \column{0.5\textwidth}
    \textbf{Direct signals:}
    \begin{itemize}
      \item \texttt{stars} — popularity proxy
      \item \texttt{forks} — adoption proxy
      \item \texttt{contributors\_count} — ecosystem breadth
      \item \texttt{weekly\_commit\_avg} — development velocity
      \item \texttt{days\_since\_last\_push} — recency / abandonment
      \item \texttt{repo\_age\_days} — maturity
      \item \texttt{releases\_count} — production-readiness
      \item \texttt{open/closed\_issues} — community activity
      \item \texttt{has\_ci} — engineering maturity
      \item \texttt{readme\_length} — documentation quality
    \end{itemize}

    \column{0.48\textwidth}
    \textbf{Derived features:}
    \begin{itemize}
      \item \texttt{activity\_score} — composite score (commits, recency, CI, issue resolution)
      \item \texttt{stars\_per\_day} — normalized growth rate
      \item \texttt{fork\_star\_ratio} — adoption vs.\ interest
      \item \texttt{issue\_resolution\_rate} — community responsiveness
    \end{itemize}

    \bigskip
    \begin{block}{Key insight}
      No single signal determines maturity. A declining repo can have
      many stars — it may have been popular and then abandoned.
    \end{block}
  \end{columns}

\end{frame}

% ============================================================
\section{LLM Weak Labeling}
% ============================================================

\begin{frame}{From signals to text summaries}

  Before labeling, each repository is converted into a \textbf{plain-English paragraph}:

  \bigskip
  \begin{exampleblock}{Example summary — \texttt{statsmodels/statsmodels}}
    \small
    \textit{``Repository `statsmodels/statsmodels' is 15.0 years old and written primarily in Python.
    It has 11,435 stars and 3,343 forks, with 457 contributors. Development velocity: very active
    with 6.2 commits per week on average, and it was updated within the last month. It has 31
    releases, 2,973 open issues, and 1 closed issue. The project has extensive README documentation
    and has CI/CD workflows configured. Overall activity score: 0.61/1.0.''}
  \end{exampleblock}

  \bigskip
  \textbf{Why natural language?}
  \begin{itemize}
    \item LLMs reason better over text than raw tabular values
    \item Verbalizing signals (``last updated 3 years ago'') encodes domain knowledge
    \item Same representation is used as input for DistilBERT fine-tuning
  \end{itemize}

\end{frame}

% ------------------------------------------------------------

\begin{frame}{DeepSeek as weak annotator}

  \begin{bluebox}{Weak supervision}
    Instead of manually labeling 587 repositories, we use \textbf{DeepSeek-V3}
    (via OpenAI-compatible API) as the initial annotator. Labels are ``weak''
    because they are machine-generated and may be noisy — the BERT model
    learns to generalize beyond individual label errors.
  \end{bluebox}

  \bigskip
  \textbf{Cost:} \$0.048 USD total for two full labeling runs (587 repos each).

  \bigskip
  \textbf{Prompt iteration — methodological sensitivity (Q4):}

  \vspace{0.3em}
  \begin{columns}
    \column{0.48\textwidth}
    \textbf{Baseline prompt:}\\
    Strict definitions.\\[0.3em]
    \small Emerging: 15 \ | \ Experimental: 122 \\
    Mature: 113 \ \ \ | \ \ Declining: 337

    \column{0.48\textwidth}
    \textbf{Refined prompt:}\\
    Nuanced + balance guidance.\\[0.3em]
    \small Emerging: 27 \ | \ Experimental: 217 \\
    Mature: 85 \ \ \ \ | \ \ Declining: 258
  \end{columns}

  \bigskip
  \small The refined prompt increased \texttt{emerging} coverage by 80\% and reduced
  \texttt{declining} dominance from 57\% to 44\% of the dataset.

\end{frame}

% ============================================================
\section{Model \& Results}
% ============================================================

\begin{frame}{DistilBERT fine-tuning pipeline}

  \begin{columns}
    \column{0.55\textwidth}
    \textbf{Architecture:} \texttt{distilbert-base-uncased}
    \begin{itemize}
      \item 66M parameters (40\% fewer than BERT)
      \item Input: repository text summary ($\leq$256 tokens)
      \item Output: 4-class softmax head
    \end{itemize}

    \bigskip
    \textbf{Training details:}
    \begin{itemize}
      \item Stratified split: 70/15/15
      \item 410 train | 88 val | 89 test
      \item 4 epochs, lr = 2e-5, batch = 16
      \item Weighted cross-entropy loss to handle class imbalance
      \item Best model at epoch 2 (Val F1 = 0.52)
    \end{itemize}

    \column{0.43\textwidth}
    \textbf{Class weights (inverse frequency):}

    \vspace{0.5em}
    \begin{tabular}{lr}
      \toprule
      Class & Weight \\
      \midrule
      \textcolor{emerging}{emerging} & 2.576$\times$ \\
      mature & 0.830$\times$ \\
      declining & 0.272$\times$ \\
      experimental & 0.322$\times$ \\
      \bottomrule
    \end{tabular}

    \bigskip
    \small Weights penalize errors on minority classes, especially \textcolor{emerging}{emerging}.
  \end{columns}

\end{frame}

% ------------------------------------------------------------

\begin{frame}{Test set results}

  \begin{columns}
    \column{0.45\textwidth}
    \textbf{Overall performance:}

    \vspace{0.5em}
    \begin{tabular}{lr}
      \toprule
      Metric & Value \\
      \midrule
      Accuracy & 55.1\% \\
      Macro F1 & 0.41 \\
      \bottomrule
    \end{tabular}

    \vspace{1em}
    \textbf{Per-class F1:}

    \vspace{0.5em}
    \begin{tabular}{lr}
      \toprule
      Class & F1 \\
      \midrule
      \textcolor{declining}{declining} & 0.64 \\
      \textcolor{mature}{mature} & 0.55 \\
      \textcolor{experimental}{experimental} & 0.47 \\
      \textcolor{emerging}{emerging} & 0.00 \\
      \bottomrule
    \end{tabular}

    \column{0.53\textwidth}
    \begin{redbox}{Error analysis}
      \begin{itemize}
        \item \textbf{Emerging (F1=0.00):} only 4 test samples — insufficient signal despite class weighting. This is a \textit{known limitation} of weak supervision with severe imbalance.
        \item \textbf{Declining vs.\ Experimental:} most frequent confusion — both share low-activity signals but differ in intentionality.
        \item \textbf{Mature:} occasionally confused with declining for very old, slow-moving but stable projects.
      \end{itemize}
    \end{redbox}
  \end{columns}

\end{frame}

% ============================================================
\section{Business Value \& Limitations}
% ============================================================

\begin{frame}{Who can use this system?}

  \begin{columns}
    \column{0.5\textwidth}
    \begin{block}{Central banks \& governments}
      Monitor which econometric tools are gaining traction to inform
      software procurement, training investment, and research infrastructure decisions.
      \textit{Ex: should a central bank adopt a new DSGE framework?}
    \end{block}

    \vspace{0.4em}

    \begin{block}{Academic institutions}
      Identify emerging methodological tools for curriculum updates and
      faculty research priorities.
    \end{block}

    \column{0.48\textwidth}
    \begin{block}{Consulting firms \& think tanks}
      Track which quantitative methods are gaining or losing traction to
      advise clients on tooling adoption.
    \end{block}

    \vspace{0.4em}

    \begin{block}{Investors in developer tooling}
      Identify growing open-source ecosystems with commercial potential
      or flag declining tools before pivoting strategies.
    \end{block}
  \end{columns}

\end{frame}

% ------------------------------------------------------------

\begin{frame}{Limitations \& future directions}

  \begin{redbox}{Key limitations}
    \begin{itemize}
      \item \textbf{Weak labels are noisy:} DeepSeek may misclassify edge cases —
            the model inherits these errors
      \item \textbf{Severe imbalance in \texttt{emerging}:} 27 samples are insufficient
            for reliable classification of this category
      \item \textbf{Static snapshot:} data reflects a single collection point —
            longitudinal signals (growth \textit{rate}) would be more informative
      \item \textbf{Selection bias:} GitHub search favors popular repos —
            niche or abandoned tools may be underrepresented
    \end{itemize}
  \end{redbox}

  \bigskip
  \textbf{Future directions:}
  \begin{itemize}
    \item Collect data at multiple time points to compute \textit{trend} signals
    \item Expand to PyPI/CRAN download statistics for adoption validation
    \item Use larger BERT models (DeBERTa-v3) with GPU training
    \item Human-in-the-loop validation of LLM labels for minority classes
  \end{itemize}

\end{frame}

% ============================================================
% CONCLUSIÓN
% ============================================================

\begin{frame}{Summary}

  \begin{greenbox}{What we built}
    A complete \textbf{weak supervision NLP pipeline} that:
    \begin{enumerate}
      \item Collects 12 GitHub signals for 587 econometrics repositories
      \item Converts signals into natural language summaries
      \item Uses DeepSeek to generate weak labels (cost: \$0.048 USD)
      \item Fine-tunes DistilBERT to classify technological maturity
      \item Exposes results via a 4-tab Streamlit application
    \end{enumerate}
  \end{greenbox}

  \bigskip
  \begin{columns}
    \column{0.5\textwidth}
    \textbf{Key takeaway:}\\[0.3em]
    The value of the system is not its accuracy (55\%) —
    it is its ability to \textbf{scale human judgment} over hundreds of
    repositories at negligible cost, providing actionable ecosystem intelligence.

    \column{0.48\textwidth}
    \textbf{Methodological contribution:}\\[0.3em]
    Prompt iteration (Q4) showed that label quality
    and class balance are more important than model architecture
    in weak supervision settings.
  \end{columns}

\end{frame}
